\subsection{The two directions a sampling representation must support}
\label{sec:threshold-positive-transfer}

Throughout this section, \emph{full AA} means classical approximation of
the acceptance probability of an arbitrary quantum query algorithm under
the specified input law. We use mean absolute error: a simulator for a
\(T\)-query acceptance function \(Q\) has
\(\mathbb E_\mu|A-Q|\le e\). Substituting \(e=\epsilon\delta\) gives
error at most \(\epsilon\) except on a \(\delta\)-fraction of inputs.
We import the uniform-input simulation consequence of AA, with query
bound \(s_U(T,e)=\operatorname{poly}(T,1/e)\), as in the existing
exchangeable-simulation section. All statements below are deductions
conditional on that imported result, not additional claims about its
provenance. Computation, preprocessing, and the storage of nonuniform
decision trees are free; only queries to the input are counted.

\begin{proposition}[A local posterior transfer criterion]
\label{prop:threshold-local-posterior}
Let \((X,Z)\) be a coupling with \(X\sim\mu\), \(Z\sim\nu\), where
\(\nu\) has a full-AA simulator costing \(s_\nu(T,e)\). Suppose a
deterministic map \(F\) has the following properties.
\begin{enumerate}
\item A coherent query to \(F(Z)_i\) can be implemented with at most
\(a\) queries to \(Z\), uniformly in \(i\).
\item For every \(x\) of positive \(\mu\)-probability and every \(i\),
\(\Pr[F(Z)_i\ne x_i\mid X=x]\le\gamma\).
\item Given oracle access to \(x\), the conditional law \(Z\mid X=x\)
can be supplied consistently on demand: answering any adaptive sequence
of \(s\) source-bit requests uses at most \(b(s)\) queries to \(x\).
\end{enumerate}
Then a \(T\)-query acceptance function has a deterministic simulator
using at most \(b(s_\nu(aT,e/2))\) queries and mean error at most
\[
 e/2+4T\sqrt\gamma.
\]
\end{proposition}

\begin{proof}
Apply simulation under \(\nu\) to the genuine quantum algorithm
\(Q\circ F\), obtaining a source decision tree \(D(Z)\) with average
error \(e/2\). On input \(X=x\), run this tree using the conditional
source oracle from the third hypothesis. Averaging over the coupling
preserves the \(e/2\) bound. The quantum hybrid argument, conditioned
on \(X=x\), gives
\(\mathbb E|Q(X)-Q(F(Z))|\le4T\sqrt\gamma\): at each query the
expected weight on disagreements is at most \(\gamma\), even though
the query weights depend on \(x\). The triangle inequality proves the
claim. The conditional oracle is an ordinary randomized query procedure
using input-independent random coins; fixing those coins by averaging
gives one deterministic simulator of no greater average error.
\end{proof}

The conditional probability in the second hypothesis is essential.
A marginal estimate \(\Pr[F(Z)_i\ne X_i]\le\gamma\) does not by
itself control quantum query weights that adapt to \(X\). Likewise,
an efficient forward sampler alone does not provide the third
hypothesis. The theorem requires sampling the hidden randomness
\emph{conditional on the actual input}, with a query bound. An
arbitrary global posterior computation might require reading all of
\(x\).

This criterion has no source-fanout requirement. One source bit may
affect many output spins, provided that its value can be sampled from
the posterior using few observed spins. That is a substantive
difference from local flip-influence transfer.
